
\documentclass[11pt]{article}
\usepackage{geometry}
\geometry{a4paper, margin=1in}
\usepackage{amsmath,amssymb}
\usepackage{booktabs}
\usepackage{graphicx}
\usepackage{microtype}
\usepackage{url}
\usepackage[hidelinks]{hyperref}
\usepackage{xcolor}
\usepackage[htt]{hyphenat}
\sloppy

\title{Post-Tuning Closed-Choice LLM Agents Across Four Model Sizes:\\
Cost Comparisons, Training Stops, and Incomplete Evaluation\\
on a Simulated Bullwhip Environment}
\author{Pepijn Frenken\\
Independent researcher\\
\texttt{pepijnfrenken@mailbox.org}}
\date{2026-10-22}

\begin{document}
\maketitle

\begin{abstract}
A published claim \cite{long2026} reports that a GRPO-tuned 4B agent beats its own untrained
baseline on cost in a four-tier beer-game environment. We report a pre-registered,
commit-dated protocol that tests a claim of that type at four model scales (0.6B, 1.7B, 4B,
8B), as an adaptation study: our training procedure and readout differ from the published
recipe in enumerated ways (Appendix~\ref{app:deltas}). The outcome vector is:
\textbf{4B passes} the frozen paired cost criteria on both evaluated seeds (medians 8{,}315
and 7{,}612 vs matched untrained control 21{,}578; 30/30 paired wins per seed, no ties;
sign-test $p = 2^{-30} = 9.3\times10^{-10}$ per seed; median-ratio reductions 61.5\%/64.7\%;
median-of-per-path-reductions 62.9\%/64.8\%; the sign test covers 30 paired evaluations of a
trained checkpoint, not 30 training seeds). \textbf{1.7B is adverse}: all three evaluated
seeds regress against their matched untrained-1.7B control (medians $1.87\times$--$6.19\times$,
0/30 wins each), and a second learning-rate arm tripped the collapse stop-loss. \textbf{0.6B}
has negative wave-7 paired comparisons on all six arm--seed rows and a wave-8 procedure
stopped by its own stop-loss (the collapse detector tripped on all six arm tasks; five were
stopped before final evaluation and one completed with a failing output-mass gate).
\textbf{8B is not evaluable within size}: no matched untrained-8B control evaluation and no
final-eval generation exist; its two learning-rate arms are reported from the
curriculum-validation band on document-level provenance (pooled means 23{,}867 vs 46{,}969;
the selected arm sits between the trained-4B band reference and the untrained-4B band
reference). The frozen seed-pooling selection rule is the methodological product: at 8B it
refuses a favorable single-seed reading that a per-seed report would have accepted. The spec,
per-run JSON, sbatch ledger, and commit hashes are cited in the text; the repository is
available on request (Section~\ref{sec:repro}).
\end{abstract}

\section{Introduction}\label{sec:intro}
Post-training LLM agents for operational decision problems is an active recipe family. The
motivating published claim for this paper is \cite{long2026}: Long, Simchi-Levi, Zhu, Su,
A.~P.~Calmon, and F.~P.~Calmon (\texttt{arXiv:2605.17036}, version 4, 21 September 2026),
which reports that a GRPO-tuned autoregressive 4B agent (Qwen-3 4B; mean total cost 1{,}585
untrained vs 952 trained, both across 30 identical MIT Beer Game runs, v4 \S5.2) beats the
untrained 4B baseline on cost in a four-tier beer-game simulator. The open question is not
whether such a claim can hold once at one scale, but whether a claim of that type holds under
pre-registered controls and whether it survives scaling. We test that on a public four-tier
beer-game style environment with our own adapted training procedure
(Section~\ref{sec:method}). In the reference study, facilities occasionally ordered enormous
multiples of demand out of the box across repeated runs, which is the reliability symptom our
closed readout first attacks \cite{aimbench}; Sterman \cite{sterman} supplies the classic
inventory-behavior background for the environment itself.

Deviations made along the way (a 0.6B run stopped by its own pre-registered stop-loss, an
unresolved learning-rate provenance question at 1.7B, an 8B budget amendment) are documented
in Section~\ref{sec:prov} rather than absorbed silently.

\paragraph{Contributions.}
\begin{itemize}
  \item A pre-registered protocol instance for post-tuning a closed-choice decision readout
  on a beer-game cost signal, with frozen pre-flight gates (grid regret F0, teacher F1,
  interface F2, control statistics F3), a fully specified collapse detector
  (Section~\ref{sec:gates}), and a frozen paired performance test; plus a chronology that
  separates planned contingencies from amendments, with commit hashes.
  \item An outcome vector across four scales (0.6B--8B) under that fixed protocol, with
  per-rung matched untrained controls where one exists, and the exact run ledger
  (Table~\ref{tab:ledger}) distinguishing recipe arms, learning-rate arms, seeds, stopped
  runs, and completed evaluations.
  \item A methodological result: the frozen seed-pooling selection rule blocks a favorable
  single-seed reading at 8B that a per-seed report would have accepted.
  \item A diagnostic observation, not a mechanism claim: on the pre-registered
  permutation-invariance check, no arm (and no control) preserved the argmax \emph{value}
  across shuffled choice grids, so the interface claim ``the model reads option values'' is
  withheld for every arm, per the pre-registered consequence clause.
\end{itemize}

\section{Environment}\label{sec:env}
The environment is a four-tier (retailer / wholesaler / distributor / factory) Beer-Game
simulator (\texttt{agent\_bullwhip/engine.py}): weekly demand, shipment delay 2 weeks,
holding cost 1.0 per unit-week, backlog cost 2.0 per unit-week, horizon 36 weeks, initial
stocks, outstanding orders, and backlogs all zero, no backlog cap or overflow penalty
(otherwise stated, nothing is truncated). Agents output one order per stage per week on a
frozen, closed ordinal choice grid of 26 letters mapping to fixed order quantities, cited
from the implemented frozen revision \texttt{make\_w8\_g26b()}, values verified from the
artifact code:
\begin{center}
0, 2, 4, 6, 8, 10, 12, 14, 16, 18, 20, 22, 24, 26, 28, 30, 36, 44, 54, 64, 78, 94, 114, 136, 166, 200.
\end{center}
Sixteen values cover the demand-mean band 0--30 in step-2 resolution; the ten tail values
are the generated geometric ladder $v_k = 2\cdot\mathrm{round}(30(200/30)^{k/10}/2)$ for
$k=1..10$ from the pre-registered revision record (\texttt{w8\_g26 -> w8\_g26b} after F0
failed on 1{,}000 probe states per regime; the failing record
\texttt{results/wave8/f0\_w8\_g26\_prereg.json} and the revision were committed before any
training, and the spec's \S3.1 pre-authorizes exactly that revision before training).
Evaluation readout is the \textbf{noisy-argmax} interface, defined now: for each decision,
the model produces logits over the 26 label tokens at the answer position; the label
distribution is $\pi( a \mid s) = \mathrm{softmax}$ over those 26 positions; the
\emph{argmax} arm reads the modal letter deterministically; the \emph{noisy-argmax} arm
samples the letter from $\pi(\cdot \mid s)$ with 30 distinct evaluation-path seeds
(1000--1029) applied at the per-decision sampling level, one simulated game per seed path.
The exact per-action expected-cost scoring used at training time is the same enumeration
without the sampling. This readout is one of the deliberate deltas from the reference
recipe's sampled-$K$ (Appendix~\ref{app:deltas}).

\section{Protocol and Gates}\label{sec:gates}
The protocol is commit-then-run: every gate and the selection rule are frozen in
version-derived spec files committed 2026-10-07 (spec commit \texttt{7d3260b}; byte-equal
cluster kit commit \texttt{4b8bffb}, same day) before any GPU job for this wave (first
training log 2026-10-08T13:19+02:00). Amendments, planned-contingency handling, and one open
item are listed in Section~\ref{sec:prov}. The frozen gates:
\begin{itemize}
  \item \textbf{F0 grid regret.} The choice grid must express the simulator's optimum in
  both demand regimes within pre-set regret and letter-mass bounds, checked on 1{,}000
  states per regime. The candidate \texttt{w8\_g26} failed F0; the revision to
  \texttt{w8\_g26b} was made via the pre-authorized pre-F0-test procedure.
  \item \textbf{F1 teacher.} The one-step-improvement teacher must beat the continuation
  floor on 30 paired curriculum paths, else no training arm is run (it passed: PI1 mean cost
  4{,}591 vs floor 8{,}541, $30/30$ paths, seeds band 9000--9029, disjoint from eval band
  1000--1029).
  \item \textbf{F2 interface.} All 26 answer letters are single tokens on the rendered
  continuation (full prompt SHA \texttt{55bdc663eee13e101be89e1ddf112e48fcafc62d4c8bf495
  0b39a2b49dafc422} for \texttt{w8\_g26b}); label ids contiguous.
  \item \textbf{F3 control statistics (collapse detector, fully specified).} The state
  \emph{answer} variable is the readout's per-state 26-letter probability row
  $\pi(\cdot \mid s)$, not a categorical state identifier (so no unique-state / deterministic
  argmax degeneracy arises); a state enters only through its row. The \emph{state-dependence}
  statistic is $I = H(\overline{\pi}) - \overline{H}$ in nats, where $\overline{\pi}$ is the
  mean answer distribution over the 1{,}024 frozen states $S_{\mathrm{nom}}$ and
  $\overline{H}$ their mean per-state entropy (\texttt{decision\_gates.py},
  \texttt{mutual\_information()}). Thresholds are frozen per model size before training from
  that model's own untrained readout on \texttt{S\_nom}: $H_{\mathrm{lo}}$ / $H_{\mathrm{hi}}$
  at the 5th/95th percentile of per-state entropy; the modal-letter share $m^*$ of the
  simulator oracle argmin letter (letter A there) and its Wilson 95\% upper bound
  $m_{\mathrm{hi}}$; and $I_{\mathrm{hi0}}$, the 95th percentile of $I$ over $B = 1{,}000$
  bootstrap resamples of the rows (rows sampled with replacement, seed 20261007). A trained
  arm is \emph{collapsed} iff [median entropy $< H_{\mathrm{lo}}$ AND (modal-argmax-letter
  share $> m_{\mathrm{hi}}$ OR $I \leq I_{\mathrm{hi0}}$)] OR [median entropy
  $> H_{\mathrm{hi}}$ AND $I \leq I_{\mathrm{hi0}}$]. Low entropy alone is not collapse.
  Frozen values: 1.7B $H_{\mathrm{lo}}{=}0.000$, $H_{\mathrm{hi}}{=}1.631$,
  $I_{\mathrm{hi0}}{=}1.486$, $m_{\mathrm{hi}}{=}0.149$; 4B
  $H_{\mathrm{lo}}{=}0.038$, $H_{\mathrm{hi}}{=}1.911$, $I_{\mathrm{hi0}}{=}0.677$,
  $m_{\mathrm{hi}}{=}0.149$; 8B $H_{\mathrm{lo}}{=}0.001$, $H_{\mathrm{hi}}{=}1.881$,
  $I_{\mathrm{hi0}}{=}1.034$, $m_{\mathrm{hi}}{=}0.149$. An untrained-8B \emph{diagnostic}
  checkpoint supplies the 8B threshold inputs; no untrained-8B final \emph{control
  evaluation} was ever run. A collapsed arm's costs are reported but cannot support the
  success claim. The stop-loss handling is a planned contingency stated in the spec
  (\S6.5), not an amendment.
  \item \textbf{Permutation invariance (diagnostic, pre-registered consequence).} 512
  states; 3 shuffled grids besides the canonical one (grid seed 20261006). A state is a
  \emph{success} iff the argmax letter's \emph{value} is unchanged on all three shuffles;
  the arm passes iff the success fraction $\geq 0.90$ of states
  (\texttt{decision\_gates.py}, \texttt{pass} field over \texttt{invariance\_rate}). The
  diagnostic's \texttt{invariance\_rate} therefore counts successes per state (all three
  shuffles simultaneously), not per state--shuffle comparison. Observed at 4B
  (\texttt{w8-r1-4b-seed0-perm.json}): rate $0.0469$ (24 of 512), per-shuffle disagreement
  counts $[0, 464, 452, 449]$ vs the canonical. Earlier arms (including the untrained
  control) read rate $0.000$. This is explicitly NOT part of the performance-claim
  conjunction: wave 7 already read the shared failure, so the pre-registered consequence
  clause says on a fail the cost gates still decide the performance claim, the interface
  claim ``the model reads option values'' is withheld, and the arm is labelled a
  \emph{fixed-ordinal-scale specialist}. We keep that label for every arm here.
  \item \textbf{G7 output-mass gate (fully specified).} For each decision, the fraction of
  full-vocabulary probability on the 26 label tokens is
  $\exp(\mathrm{logsumexp}(\mathrm{label\ logits}) - \mathrm{logsumexp}(\mathrm{full\ vocab\ logits}))$
  (\texttt{decision\_readout.py}); a run's \texttt{min\_mass} is the minimum of that
  per-decision mass over the run; a summary's \texttt{min\_mass} is the minimum over runs.
  G7 for a trained arm: summary \texttt{min\_mass} $\geq 0.9$ (i.e. the smallest
  min-over-decisions mass across the arm's runs must still leave 90\% of probability on
  the letter set). It is possible (an individual letter probability cannot exceed 1/26, but
  the \emph{total} mass on the 26 letters can fall to 0, and a readout that has drifted off
  the label set scores far below 0.9). G7 is mandatory for a trained arm's eligibility
  (from wave-7 onwards, where observed values 0.925--0.983 motivated the 0.9 threshold).
  Results for the passing 4B checkpoints: seed 0 \texttt{min\_mass} $0.99999$, seed 1
  $0.99999$ (both PASS). The 1.7B seed-1 arm's $7.3\times10^{-9}$ is a binding G7
  \emph{failure} on a negative arm; it does not change that arm's verdict but is recorded.
  \item \textbf{Selection rule (per learning-rate arm, frozen).} Within each pre-registered
  learning-rate arm, seeds are pooled by \emph{mean} validation cost on the fixed
  curriculum-validation band; the arm with the lower pooled mean is selected and the other
  is reported as sensitivity. The selected arm must then beat its rung's matched untrained
  control under the frozen paired test G1: median cost reduction $\geq 15\%$ AND
  median-of-per-path-reductions $\geq 15\%$, $\geq 20/30$ paired wins, one-sided exact sign
  test $p < 0.02$; a seed passes iff all three hold. The criteria have a frozen redundancy:
  at $n=30$ non-tied pairs, 20 wins gives $p_{\mathrm{one-sided}} = 0.049$, 21 gives 0.021,
  and 22 gives 0.008, so the $p<0.02$ clause implies $\geq 22/30$ wins. Tie handling: a tie
  is neither win nor loss; all 60 paired comparisons here had zero ties (verified from the
  raw JSONLs), and rows with a zero control cost would be dropped from the denominator
  (none occurred). \textbf{The pass rule ``the arm needs $\geq 2/3$ passing seeds'' is a
  fraction of planned seeds} (\texttt{WAVE8\_SPEC.md} \S6.3, committed 2026-10-07, i.e.
  before any wave-8 run): three planned seeds at 1.7B ($\geq$ 2 of 3 must pass), two planned
  seeds per learning-rate arm at 4B and 8B (both evaluated seeds must pass). A planned seed
  that never becomes an evaluated checkpoint does not enter the fraction; it is reported as
  missing (1.7B's 5e-5 arm: one of three planned seeds trained, stopped at round 2, zero
  evaluated).
\end{itemize}

\section{Training objective and configuration}\label{sec:method}
We deliberately avoid borrowing algorithm names for what we run. The trained objective at
every rung is, per state $s_i$ and over the frozen 26-letter choice set $A$:
\begin{align}
L(\theta) &= \frac{1}{N}\sum_{i=1}^{N}\bigg[\sum_{a \in A} \pi_\theta(a \mid s_i)\, z_i(a)
+ \beta\, D_{\mathrm{KL}}\!\left(\pi_\theta(\cdot \mid s_i)\,\Vert\,\pi_{\mathrm{ref}}(\cdot \mid s_i)\right)\bigg],\\
z_i(a) &= \frac{J(s_i,a)-\overline{J}_i}{\max(\mathrm{std}_{a}J_i,\,10^{-3})}.
\end{align}
Here $J(s_i,a)$ is the simulator's expected total cost of playing action $a$ from state
$s_i$ and then continuing with the frozen reference policy for the remainder of the game,
averaged over the labeler's continuation rollouts; $\overline{J}_i$ and $\mathrm{std}_a J_i$
are per-state statistics over the full choice set; $\pi_{\mathrm{ref}}$ is a frozen
policy-independent anchor. The KL term is averaged inside the state sum (one KL per state,
then averaged over the $N$ states), $\beta = 0.02$ in all wave-8 runs
(\texttt{state.json}, per-round training logs). Per round: fresh collection over the
curriculum band (argmax and sampled passes re-labelled each round, 5{,}168 training / 272
validation states at 4B), 162 batches per round, batch size 32, optimizer AdamW, learning
rate per arm (0.6B: 1e-5 and 5e-5; 1.7B: 5e-5 planned, 1e-5 executed; 4B: 5e-5; 8B:
5e-5 and 1e-5), LoRA $r{=}8$, $\alpha{=}16$, dropout 0.05, attention-only targets
(\texttt{q\_proj, k\_proj, v\_proj, o\_proj}), checkpoint at step 169 selected by the
frozen validate stage. Three rounds per arm. The gradient is \emph{exact} in a qualified
sense: the action sum enumerates every member of the 26-value choice set, and $J$ is
deterministic per (state, action) at scoring time; simulator-level stochasticity (demand
draws inside $J$'s continuation) is part of $J$'s value, not resampled per gradient step.
Deltas from the reference GRPO objective \cite{grpo,long2026} are enumerated in
Appendix~\ref{app:deltas}; notably the reference objective (v4 eq.~2) also carries a
frozen-KL term to its reference policy, which this recipe shares, and the reference
normalizes advantages per sampled group while this recipe $z$-scores per state over the
enumerated set. Published RLCD \cite{rlcd} builds contrastive preference pairs and a
preference model; no preference model exists here. Both are cited as lineage, not as what
we execute; a verbatim re-run of either is future work.

\section{Run ledger}\label{sec:ledger}
Table~\ref{tab:ledger} is one row per planned task (recipe arm $\times$ learning-rate arm
$\times$ seed), derived from the per-run artifacts; all subsequent counts come from this
table and nowhere else.

\begin{table}[ht]
\centering
\caption{Planned-task ledger. ``Stopped R$k$'' = collapse stop-loss persisted at round
$k$; ``eval'd'' = matched noisy-argmax final evaluation exists.}
\label{tab:ledger}
\small
\begin{tabular}{@{}p{2.2cm}p{2.6cm}p{1.6cm}p{5.4cm}p{1.4cm}@{}}
\toprule
\textbf{Rung} & \textbf{Recipe lr} & \textbf{Seeds} & \textbf{Task status} & \textbf{Eval'd} \\
\midrule
0.6B w7 & SFT & 0,1,2 & 3 arm--seed rows, completed C1 protocol & 3 \\
0.6B w7 & exact-grad & 0,1,2 & 3 arm--seed rows, completed C1 protocol & 3 \\
0.6B w8 R1 & 1e-5 & 0 & stopped R2, accidentally resumed to R3 (markers both DONE/STOPPED) & 0 \\
0.6B w8 R1 & 1e-5 & 1,2 & stopped R2 & 0 \\
0.6B w8 R1 & 5e-5 & 0 & stopped R2 (round reads C, clean, C) & 0 \\
0.6B w8 R1 & 5e-5 & 1,2 & stopped R2 & 0 \\
1.7B & 5e-5 & 0 & stopped R2 & 0 \\
1.7B & 5e-5 & 1,2 & \multicolumn{2}{l}{not attempted (stop preceded further 5e-5 spend)} \\
1.7B & 1e-5 & 0,1,2 & 3 arm--seed rows, completed R3 & 3 \\
4B & 5e-5 & 0,1 & 2 arm--seed rows, completed R3 & 2 \\
8B & 5e-5 & 0,1 & trained on external lanes; validation band only (23{,}867 pooled) & 0 \\
8B & 1e-5 & 0,1 & trained on external lanes; validation band only (46{,}969 pooled) & 0 \\
\bottomrule
\end{tabular}
\end{table}

Counts derived from Table~\ref{tab:ledger}: 33 planned tasks across 12
recipe$\times$lr$\times$seed expansions; 21 tasks reached a trained checkpoint; 11 rows
have a matched paired evaluation (0.6B w7 6, 1.7B 3, 4B 2); 8 arm--seed rows trained but
carry no final evaluation (0.6B w8 six, 1.7B 5e-5 one, and the 8B validation band rows
separately). Wave-7's ``all six arm--seed rows'' in the results text counts rows of the
two-recipe $\times$ three-seed structure above, not six independent recipe arms.

\section{Results}\label{sec:results}
Table~\ref{tab:ladder} is the headline table. Lower cost is better. Every row names its own
estimand explicitly: the statistic, the evaluation generation, and the control it is
compared to. Paired wins are per-path (paired on the demand path index), one-sided sign
test.

\begin{table}[ht]
\centering
\caption{Outcome vector. ``final eval'' = matching noisy-argmax $n{=}30$ generation from
\texttt{eval/} (same adapters as the \texttt{wave8\_val} selection band); medians on the
first line, means on the second where used; ``band'' = curriculum-validation means.}
\label{tab:ladder}
\small
\begin{tabular}{@{}p{1.5cm}p{2.5cm}p{4.4cm}p{3.9cm}p{1.9cm}@{}}
\toprule
\textbf{Rung} & \textbf{Arms} & \textbf{Estimand} & \textbf{Values vs control} & \textbf{Status} \\
\midrule
0.6B w7 & SFT + exact-grad & C1, noisy-argmax n=30, per arm--seed & 0/30 wins on all six arm--seed rows & adverse \\
0.6B w8 & exact-grad & stopped before final eval & collapse tripped on all 6 tasks; 5 stopped pre-eval & stopped \\
1.7B & exact-grad, 3 seeds & final-eval medians & ctl 11{,}508; medians 21{,}468/43{,}660/71{,}270 & adverse (0/3) \\
4B & exact-grad, 2 seeds & final-eval medians and means & ctl 21{,}578; \textbf{pass both seeds} & \textbf{pass (2/2)} \\
8B & exact-grad, 4 legs & validation-band means only & no matched untrained-8B final control & \textbf{not evaluable} \\
\bottomrule
\end{tabular}
\end{table}

\paragraph{4B, the passing rung, with both effect sizes and both statistics.} Both seeds
pass the frozen G1 in full on the final-eval generation vs the matched untrained-4B
control:
\begin{itemize}
  \item median cost: 8{,}315 (seed 0) and 7{,}612 (seed 1) vs control 21{,}578; mean cost
  8{,}286.9 and 7{,}585.3 vs control mean 27{,}886 (\emph{recomputed from the raw paired
  JSONLs for this section}, the first time both statistics appear together);
  \item median-ratio reduction: $1 - 8{,}315/21{,}578 = 61.5\%$ and
  $1 - 7{,}612/21{,}578 = 64.7\%$ (criterion $\geq 15\%$);
  \item \emph{median of the 30 per-path proportional reductions}
  ($\mathrm{median}_i(1 - c^{\mathrm{train}}_i / c^{\mathrm{ctl}}_i)$, zero-control rows
  dropped; zero such rows, zero ties): 62.9\% and 64.8\% (recomputed from the raw pairs);
  \item paired wins: 30/30 per seed, one-sided exact sign test
  $p = 2^{-30} = 9.3\times10^{-10}$ per seed (criterion effectively $\geq 22$ wins from
  the $p<0.02$ clause);
  \item G7 output-mass gate: seed 0 \texttt{min\_mass} 0.99999, seed 1 0.99999 (PASS);
  \item G2 band tier (pre-registered report tier, not a gate): noisy walk-forward tuned
  benchmark floor 5{,}346 from \texttt{results/wave7\_floors.json}; tiers are nested flags,
  applied as \emph{floor-parity} $\leq 1.25\times$ (implies \emph{headline} $\leq 2\times$,
  implies \emph{strong} $\leq 4\times$); a run receives all tiers it falls in. Both 4B
  seeds sit at $1.42\times$--$1.56\times$, i.e.\ headline and strong, not floor-parity.
\end{itemize}
Inference scope of the sign test: 30 paired evaluations of a trained checkpoint against
the frozen control on the fixed demand paths; not evidence from 30 training seeds, not a
test of a scaling claim. The two seed-level tests share the same control paths, so they
are dependent reads, reported separately and never combined.

\paragraph{1.7B, a within-size regression with a matched control.} Final-eval medians
21{,}468 / 43{,}660 / 71{,}270 (seeds 0/1/2 of the executed 1e-5 arm) against the matched
untrained-1.7B control of 11{,}508 (mean 11{,}405): 0/30 paired wins on every seed, median
regressions $1.87\times$, $3.79\times$, $6.19\times$. The curriculum-validation generation
(means 47{,}125 / 36{,}734 / 61{,}073; pooled 48{,}311) is the frozen selection input; both
generations agree on the outcome. Training made this rung worse than its own starting
readout. Separately, the planned 5e-5 arm trained seed 0 only and tripped the collapse
gate persistently (round-1 modal-argmax share 1.0 vs bound 0.149), stopping at round 2;
seeds 1 and 2 were not attempted. That arm has no validation and no evaluation, so it
neither strengthens nor weakens the 1.7B adverse result; it is an incomplete planned arm.

\paragraph{0.6B, two training variants and a stopped procedure.} Wave 7 ran two recipe
arms (SFT and exact-gradient R1) at three matched seeds each against the matched
untrained-0.6B control (median 11{,}569): the recomputation in
\texttt{VERIFY-W7-EVAL.md} matches 0/30 on every arm--seed row of the binding
noisy-argmax read (arm medians 20{,}594 / 25{,}721 / 25{,}508 SFT and 37{,}384 /
32{,}516 / 27{,}652 exact-grad; exact-gradient rows historically labelled ``RLCD'' in
internal files; sign $p = 1.0$ on each row; the exact-grad arms also showed letter-lock,
1--2 distinct order values). Wave 8's exact-gradient instance at the same scale tripped
the collapse gate on six of six arm tasks at least once (round-1 modal-argmax share
0.51--1.00 vs the bound 0.149); five tasks were stopped before final evaluation. The
sixth (lr 1e-5, seed 0) was stopped at round 2, then accidentally resumed by an
unattended bare-array resubmission and completed round 3 with markers \texttt{DONE} and
\texttt{STOPPED\_collapse} both present; its final read is not collapsed (energy
1.604 nats) but fails G7 (\texttt{min\_mass} $3.2\times10^{-4}$, the readout left the
label set). No wave-8 0.6B arm has a matched final evaluation, so the wave-8 0.6B
outcome is a stopped procedure, distinct from a measured cost regression; the wave-7
rows carry the measured negative.

\paragraph{8B, validation-band evidence only, within-size not evaluable.} All four 8B
planned legs (2 lr $\times$ 2 seeds) trained on external A100 lanes (Snellius jobs
\texttt{28414248} and \texttt{28418615} for the seed-0 legs; two Modal legs completed the
seed-1 legs after the budget amendment, Section~\ref{sec:prov}). Document-recorded
validation-band means (the declared selection input, $n{=}30$ argmax runs on the
curriculum band): 5e-5 arm 35{,}289 / 12{,}444 (pooled 23{,}867); 1e-5 arm 71{,}855 /
22{,}083 (pooled 46{,}969). We state what these numbers can and cannot support:
\begin{itemize}
  \item \emph{Selection}: the pooled band means select the 5e-5 arm under the frozen rule
  (23{,}867 $<$ 46{,}969). The selection record is document-level (verdict doc with
  leg-by-leg timings); no selection JSON was recovered from either lane, so selection
  provenance is incomplete, and the record is presented as such.
  \item \emph{Band-level, matched comparator}: on the same band, the trained-4B
  reference is pooled mean 12{,}614.8 (10{,}547.8 / 14{,}681.8, archived) and the
  untrained-4B curriculum-band mean is 116{,}948 (archived). The selected 8B arm's
  23{,}867 is therefore $1.89\times$ the trained-4B band reference and $0.20\times$ the
  untrained-4B band reference. These are band-level comparisons; they are consistent
  with the reference paper's own band structure, but they are not the final-eval
  generation, and no matched untrained-8B \emph{final control} exists, so neither
  comparison establishes an 8B within-size verdict.
  \item \emph{Seed spread inside the selected arm}: 35{,}289 vs 12{,}444 (ratio
  $2.84\times$); the other arm's ratio is $3.25\times$ (larger). The frozen
  seed-pooling rule takes the arm mean, not the favorable seed.
  \item \emph{Within-size G1}: not computable. No matched untrained-8B control
  \emph{evaluation} exists (the frozen F3 threshold artifact is a diagnostic checkpoint,
  a different object measured on a different state sample), and no final-eval generation
  was produced at 8B at all. The 8B row's within-size verdict is genuinely unevaluated,
  not negative, and no cross-size cost-regression claim is made from it.
\end{itemize}

\paragraph{Permutation diagnostic; what it does and does not show.} The pre-registered
consequence of the shared near-zero permutation-invariance read across trained arms and
controls (0.000 on the untrained control and every wave-7 arm; 0.047/0.039 on the passing
4B seeds, both below the 0.90 gate) is a label and a withheld interface claim. Zero
observed successes is not proof that the underlying success probability is zero. Our
stated \emph{hypothesis} (proposed, not demonstrated) is that the passing 4B arm is a
specialist scorer on the trained choice set: a fixed mapping from visited states to fixed
output letters that happens to order well on the eval distribution. The shared diagnostic
does not separate the passing rung from failing rungs and therefore does not by itself
explain the between-size outcome difference; no intervention isolates a mediator, so we
offer no causal account of any between-size difference.

\paragraph{Interface validity, precisely stated.} A 26-action constrained letter-logit
readout structurally guarantees that the selected action is in the set; a free-form lane
does not. On the program's own noisy-argmax generations, wave-7 and wave-8 records show
zero \texttt{parse\_fail} at the decision level: across 14 archived arm--generation
summary files and 60{,}480 decisions (30 runs $\times$ 144 decisions per run $\times$
14 generations in the paired-comparison set), there are 0 invalid-decision events, and no
run contains one. Wave-6 free-form arms under that lane's explicit failure-handling rule
(a free-form arm's out-of-contract order is dropped and flagged) carried 7--188 such
events per arm. A paired free-form-vs-closed-choice cost test at matched conditions was
not run, so the recorded effect is the \emph{elimination of invalid actions and the
simplification of the comparison}, not a demonstrated total-cost advantage.

\section{Provenance and amendments}\label{sec:prov}
We separate planned contingencies (handling pre-written in the spec before results),
actual amendments (documented at the time), and open items:
\begin{itemize}
  \item \textbf{Frozen before wave-8 training.} Spec and wave-8 pre-registration committed
  2026-10-07, sha \texttt{7d3260b} (\texttt{docs/waves/WAVE8\_SPEC.md}, with the failing
  grid record \texttt{results/wave8/f0\_w8\_g26\_prereg.json}, the revised grid
  \texttt{w8\_g26b}, the F0/F1/F2/F3 and G1/G2/G7 definitions, and the seed-pooling
  selection rule); byte-equal cluster kit commit \texttt{4b8bffb}, same day. First wave-8
  GPU job: 0.6B R1 training log starts 2026-10-08T13:19 (+02:00), i.e.\ after both
  commits. Wave-7 observations predate this and explicitly informed the permutation
  consequence clause, whose text was frozen in the wave-8 spec before wave-8 training;
  it is a planned contingency of that spec, not an amendment.
  \item \textbf{Planned contingencies, not amendments.} (a) The collapse stop-loss
  (\S6.5) and the marker semantics (terminal; a resume past a stop needs a written
  override) were frozen before training. (b) The G1 pass rule, seed-pooling selection
  rule, and all thresholds in Section~\ref{sec:gates} were frozen before training.
  \item \textbf{Amendments actually made (documented at the time).} (a) 8B budget
  amendment: two seed-1 legs were re-run on a second lane (Modal A100) after the first
  lane's GPU budget ended; per-rung thresholds stayed the frozen
  \texttt{gates\_8b.json}. Nominal effective batch 32 was matched across lanes via a
  grad-accum split (micro-batch 8 vs 4); we do not claim optimizer-step-for-step or
  numerical equivalence across providers, and no equivalence test was run.
  (b) 4B ran a single lr arm (5e-5) rather than both: the operating plan stages each
  rung at the previous rung's selected lr; the 4B result below is reported for the
  evaluated 5e-5 checkpoints as they are, and completeness of rate selection at 4B is
  treated as a separate design question, not as support for the pass.
  \item \textbf{Open provenance item (stated as open).} The 1.7B operating plan
  (\texttt{W8-START-17B.md}, \texttt{W8-NEXT-SUBMITS.md}) names 5e-5 as the planned lr;
  executed artifacts show 5e-5 seed 0 trained and collapsed to a stop at round 2 with
  no final evaluation, while the seeds that reached final eval ran at 1e-5. No amendment
  note for that switch was found in the committed docs. Regrettably the lr used for the
  \emph{evaluated} 1.7B arm is therefore not the lr the plan document names, and the
  executed-1e-5 evidence does not by itself constitute a protocol-compliant 5e-5 1.7B
  read. The 1.7B adverse result stands on the executed arm's own numbers; the plan-vs-
  execution discrepancy is carried forward unresolved rather than reconstructed from
  submit history.
  \item \textbf{Result-inspection ordering.} The curriculum-validation band
  (\texttt{wave8\_val}) was the selection input; the final-eval generation
  (\texttt{eval/}, newer by 17 h at 4B and 31 h at 1.7B, same adapters at
  checkpoint-169) was produced afterwards. When final-eval numbers were first inspected,
  the selection rule had already fixed the arm choice, so no selection step used the
  final-eval generation; per-seed report cuts made after inspection are not used to
  change any verdict.
\end{itemize}

\section{Analysis}\label{sec:analysis}
\paragraph{What the evidence supports.} One rung (4B) passes the frozen cost criteria on
both evaluated seeds, with both effect sizes, both central statistics, and the G7 gate
all recomputed from raw pairs. One rung (1.7B) is adverse on all three evaluated seeds
against a matched control, with an incomplete planned second learning-rate arm. One rung
(0.6B) is negative on all six wave-7 arm--seed rows, and its wave-8 exact-gradient
procedure is a stopped run (collapse tripped on all six tasks; five stopped before
evaluation; one completed but failed the output-mass gate). One rung (8B) is not
evaluable within size, and its band-level numbers are reported with document-level
provenance only. One-line summary: closed-choice readouts eliminated invalid actions and
simplified the comparison; exact-enumeration outcome-scored post-training paid off at
exactly one scale point (4B) under the frozen rule, and the other three scale points are
adverse, stopped, or unevaluated respectively.

\paragraph{The selection rule as a methodological point.} The 8B seed spread inside the
selected arm (35{,}289 vs 12{,}444; $2.84\times$) is the worked example. A per-seed
report can pick the flattering seed; the frozen in-arm pooled mean cannot. The
aggregation rule was frozen in the spec before outcomes were reviewed by the analyst
(commit-pinned, dated); the realized pooled mean was computed by a non-blinded analyst
when the validation outputs landed. The rule removes post-hoc interpretive freedom, not
analyst hindsight; we make no blinding claim.

\section{Limitations}\label{sec:limits}
\begin{itemize}
  \item Simulated environment; external validity is a shared problem of this study and
  the reference claim.
  \item \textbf{Bookkeeping, exact.} From Table~\ref{tab:ledger}: 33 planned tasks, 21
  trained checkpoints, 11 rows with matched paired evaluations, 13 rows stopped or
  unevaluated (6 wave-8 0.6B, 2 unevaluated 1.7B 5e-5 seeds, 1 stopped 1.7B 5e-5 seed-0
  checkpoint, 4 8B validation-band rows). Every within-size verdict is conditional on its
  pre-registered seed count; with this many experimental units there is no statistical
  test of any scaling claim, and none is offered. Cross-size comparisons are descriptive.
  \item The training objective is an exact-enumeration policy gradient with a frozen
  per-state KL anchor; the reference objective also uses a frozen reference KL (v4
  eq.~2), so that component is shared, not a delta (Appendix~\ref{app:deltas}). A
  verbatim re-run of the reference GRPO recipe or of published RLCD \cite{rlcd} is
  future work, and this study's positive rung is conditional on our recipe deltas.
  \item Hardware differs across rungs (L4 for 0.6B--4B; A100 lanes for 8B), and the 8B
  run spanned two providers; the nominal effective batch was matched but optimizer-step
  and numerical equivalence across providers were not established.
  \item No matched untrained-8B control evaluation and no 8B final-eval generation
  exist; the 8B row's within-size verdict is genuinely unevaluated, not negative.
  \item The 1.7B learning-rate provenance item in Section~\ref{sec:prov} is unresolved,
  and it further narrows what the 1.7B rows can claim protocol-compliance for.
  \item No bullwhip-amplification statistic, order-trajectory plot, or reliability-tail
  analysis is supplied; costs alone do not measure the reliability dimension the
  reference paper targets (its headline is variance compression).
\end{itemize}

\section{Reproducibility statement}\label{sec:repro}
The audit trail is tiered, stated as such. (a) Full per-run JSON on the primary cluster
archive: 0.6B wave-7 arms (6 rows, all three readout modes + permutation JSON),
untrained 0.6B/1.7B/4B, trained 1.7B (both generations), trained 4B (both seeds,
permutation JSON included), and the frozen per-rung gate files
(\texttt{gates\_1.7b/4b/8b.json}) with per-rung model revisions
(\texttt{config.sh}: Qwen/Qwen3-0.6B \texttt{c1899de2...}, Qwen3-1.7B
\texttt{70d244cc...}, Qwen3-4B \texttt{1cfa9a72...}, Qwen3-8B \texttt{b968826d...}).
(b) Present but incomplete: the wave-8 0.6B R1 arms (round structure, collapse reads,
training logs, stop markers, and the accidental-resume record present; their
final-eval files never existed, because the run was stopped before evaluation -- this
is a stop, not an archive defect) and the 1.7B 5e-5 arm (same). (c) For 8B,
document-level provenance only (verdict doc with leg-by-leg timings, sbatch/Modal job
IDs, lane metadata); no 8B raw JSONLs of either generation were recovered. Simulator
and trainer module paths: \texttt{agent\_bullwhip/engine.py} (simulator),
\texttt{decision\_eval.py} (evaluation + readout + \texttt{min\_mass}),
\texttt{decision\_train.py} (objective, LoRA, rounds, optimizer),
\texttt{decision\_readout.py} (grid registry + label mass invariant),
\texttt{decision\_gates.py} (collapse detector, permutation probe, G7); all under the
repository root cited in the ledger. Library stack declared per provider: torch 2.1.2 /
transformers 4.55.4 / peft 0.21.1 as installed by the staged sbatch kit; the runtime
manifest was not archived per lane, so these are declared versions, not runtime-verified
ones. Raw pairs regenerate every number in the results section via
\texttt{scripts/g1\_compute.py}-equivalent scripts referenced in the ledger. Numbers in
this paper are labeled recomputed-from-raw or document-recorded per row; recovering raw
8B evals would be new evidence and would not silently replace the historical rows. The
repository and artifacts are available on request.

\section*{Acknowledgments}
This is an unfunded independent replication program. The authors of \cite{long2026} are
the reference pole, not co-authors. Thanks to the TU/e Umbrella compute staff for
generous L4 time.

\bibliographystyle{plain}
\begin{thebibliography}{9}\itemsep 0pt
\bibitem{long2026} C.~X. Long, D.~Simchi-Levi, F.~Zhu, H.~Su, A.~P.~Calmon,
F.~P.~Calmon, ``Reliability and Effectiveness of Autonomous AI Agents in Supply Chain
Management,'' \texttt{arXiv:2605.17036}, version 4 (21 September 2026); v1 2026-05-16;
accessed 22 October 2026.
\bibitem{rlcd} K.~Yang, D.~Klein, A.~Celikyilmaz, N.~Peng, Y.~Tian,
``RLCD: Reinforcement Learning from Contrastive Distillation for Language Model
Alignment,'' \texttt{arXiv:2307.12950} (2023; ICLR 2024).
\bibitem{grpo} Z.~Shao, P.~Wang, Q.~Zhu, R.~Xu, J.~Song, X.~Bi, H.~Zhang, M.~Zhang,
Y.~K.~Li, Y.~Wu, D.~Guo, ``DeepSeekMath: Pushing the Limits of Mathematical Reasoning
in Open Language Models,'' \texttt{arXiv:2402.03300} (2024). (GRPO.)
\bibitem{lora} E.~J.~Hu, Y.~Shen, P.~Wallis, Z.~Allen-Zhu, Y.~Li, S.~Wang, L.~Wang,
W.~Chen, ``LoRA: Low-Rank Adaptation of Large Language Models,'' \texttt{arXiv:2106.09685}
(2021).
\bibitem{aimbench} X.~Zhao, Y.~Xie, C.~Chen, Y.~Sun, ``AIM-Bench: Evaluating
Decision-making Biases of Agentic LLM as Inventory Manager,'' \texttt{arXiv:2508.11416}
(2025).
\bibitem{sterman} J.~D. Sterman, ``Modeling Managerial Behavior: Misperceptions of
Feedback in a Dynamic Decision Making Experiment,'' \textit{Management Science} 35(3),
321--339 (1989).
\end{thebibliography}

\appendix
\section{Recipe $\times$ protocol delta table}\label{app:deltas}
Rows are split by component; the KL-reference row is shared (the reference's v4
objective eq.~2 already penalizes $D_{\mathrm{KL}}(\pi_\theta \Vert
\pi_{\mathrm{ref}})$ to a frozen $\pi_{\mathrm{ref}}$), so it is NOT a delta.
\begin{table}[ht]
\centering
\caption{\cite{long2026} v4 vs this paper, per component.}
\begin{tabular}{@{}p{3.4cm}p{4.6cm}p{4.6cm}@{}}
\toprule
\textbf{Component} & \textbf{Reference (2605.17036 v4)} & \textbf{This paper} \\
\midrule
Cost/advantage normalization & group-sampled form over $G$ rollouts & per-state
$z$-score enumerated over all 26 choices \\
KL reference policy & frozen $\pi_{\mathrm{ref}}$ (v4 eq.~2), shared & frozen
per-state anchor (shared) \\
Readout & free-form autoregressive orders & closed-choice letter logits (26 letters) \\
Reward attribution & rollout (return-to-go) vs episode & exact per-(state,action)
continuation cost $J$ \\
Model & one shared backbone, four echelons (global); per-echelon models
(decentralized) & one shared backbone across roles \\
Reward aggregation studied & agent-level vs system-level & system-level continuation
cost (closest to rollout + system scope) \\
Training-run seed count & not stated in the v4 text we could read & pre-registered per
rung (Table~\ref{tab:ledger}) \\
\bottomrule
\end{tabular}
\end{table}

\section{Artifact ledger excerpts}\label{app:ledger}
Eval generations: \texttt{eval/w7-*} (wave-7 C1 arms), \texttt{eval/w8-untrained-*}
(untrained controls), \texttt{eval/w8-r1-1.7b-seed\{n\}-*} and
\texttt{eval/w8-r1-4b-seed\{n\}-*} (final-eval generation for the rungs that produced
one); \texttt{wave8\_val/} is the curriculum-validation band used by the frozen
selection rule. Key frozen gate artifacts: \texttt{gates\_1.7b.json},
\texttt{gates\_4b.json}, \texttt{gates\_8b.json} (per-rung collapse thresholds),
\texttt{results/wave7\_floors.json} (G2 reference floors). Commit anchors:
\texttt{7d3260b} (2026-10-07, spec), \texttt{4b8bffb} (2026-10-07, byte-equal cluster
kit), first 0.6B R1 training log 2026-10-08. External 8B lanes: Snellius \texttt{28414248}
(5e-5 seed 0) and \texttt{28418615} (1e-5 seed 0), and two Modal A100 legs for the
seed-1 legs (IDs in the verdict doc). Full per-run JSON archive for the 0.6B w7 rows,
the untrained controls, and the 1.7B and 4B final generations; the wave-8 0.6B and 1.7B
5e-5-interrupted arms were stopped before evaluation (stop, not archive defect), and
the 8B lanes produced document-level records only (no recovered raw JSONLs).
\end{document}
